% ============================================================
\setpartpreamble{\partintro
  {Schema-Änderungen erzeugen, die Werkzeuge von dotnet ef erweitern und Compiled Models
  unterstützen.}
  {
    \item Kapitel~\ref{ch:how-migrations-work}: Wie Migrations funktionieren
    \item Kapitel~\ref{ch:declarative-sql-from-annotations}: Deklaratives SQL aus
    Annotations
    \item Kapitel~\ref{ch:custom-migration-operations}: Eigene Migrations-Operationen
    \item Kapitel~\ref{ch:design-time-services}: Design-Time-Services
    \item Kapitel~\ref{ch:compiled-models-trimming-and-nativeaot}: Compiled Models, Trimming
    und NativeAOT
    \item Kapitel~\ref{ch:scaffolding-reverse-engineering}: Scaffolding (Reverse
    Engineering)
  }}
\part{Migrations und Design-Time}\label{part:VI}

% ------------------------------------------------------------
\chapter{Wie Migrations funktionieren}\label{ch:how-migrations-work}

Migrations vergleichen zwei Modelle und erzeugen daraus Operationen und SQL. Ein Plugin kann
sich an mehreren Stellen einklinken; um die richtige zu wählen, muss man den Ablauf kennen:

\begin{flowdiagram}
  \flowstep{Design-Time-Modell}{Modell inklusive aller Design-Time-Annotations}
  \flowstep{IRelationalAnnotationProvider}{Überträgt Annotations auf Tabellen, Spalten,
  Indizes des relationalen Modells}
  \flowstep{IMigrationsModelDiffer}{Vergleicht altes (Snapshot) und neues relationales
  Modell}
  \flowstep{MigrationOperations}{CreateTable, AlterTable, ... mit .Annotation /
  .OldAnnotation}
  \flowstep{ICSharpMigrationsGenerator}{Schreibt Migration, Designer-Datei und Snapshot
  (dotnet ef)}
  \flowstep{IMigrationsSqlGenerator}{Erzeugt SQL beim Anwenden oder Skripten}
\end{flowdiagram}

Für Plugins gibt es zwei grundsätzliche Strategien. \textbf{Deklarativ}: Das Plugin
beschreibt den gewünschten Zustand als Annotations; EF erkennt Änderungen von selbst, und
das Plugin muss nur SQL erzeugen. \textbf{Imperativ}: Das Plugin bietet eigene Operationen
an, die Nutzer von Hand in Migrationen schreiben. Gute Plugins nutzen beide: deklarativ für
alles, was zum Modell gehört, imperativ für alles andere.

\begin{note}
  \code{Database.\allowbreak{}Ensure\allowbreak{}Created(\allowbreak{})} nutzt dieselbe
  Pipeline: Es berechnet den Diff vom leeren Zustand zum Design-Time-Modell und führt das
  SQL des Migrations-Generators aus. Eigene Migrations-SQL eines Plugins wirkt also auch
  dort, was Integrationstests erheblich vereinfacht (Kapitel~\ref{ch:test-strategy}).
\end{note}

% ------------------------------------------------------------
\chapter{Deklaratives SQL aus Annotations}\label{ch:declarative-sql-from-annotations}

\section{Der Annotation-Provider}

Der Differ vergleicht das \emph{relationale} Modell (Tabellen, Spalten), nicht die
Entity-Typen. Welche Annotations eine Tabelle trägt, entscheidet der
\code{IRelational\allowbreak{}Annotation\allowbreak{}Provider}. Den gibt es pro Context
einmal, und er gehört dem Provider; das Plugin erbt deshalb von der Npgsql-Implementierung:

% source: src/Acme.EntityFrameworkCore.Tenancy.PostgreSQL/Metadata/Internal/TenancyAnnotationProvider.cs
\begin{sourcecode}[Metadata/\allowbreak{}Internal/\allowbreak{}Tenancy\allowbreak{}Annotation\allowbreak{}Provider.\allowbreak{}cs]{csharp}
public class TenancyAnnotationProvider : NpgsqlAnnotationProvider
{
    public TenancyAnnotationProvider(RelationalAnnotationProviderDependencies dependencies)
        : base(dependencies)
    {
    }

    public override IEnumerable<IAnnotation> For(ITable table, bool designTime)
    {
        foreach (var annotation in base.For(table, designTime))
            yield return annotation;

        if (!designTime)          // nur für Migrations: Laufzeitmodell bleibt schlank
            yield break;

        // TPH/Table Splitting: der Typ mit RLS-Konfiguration, nicht einfach der erste
        var entityType = table.EntityTypeMappings
            .Select(m => (IEntityType)m.TypeBase)
            .FirstOrDefault(e => e.HasRlsConfiguration())
            ?? (IEntityType)table.EntityTypeMappings.First().TypeBase;
        foreach (var annotation in entityType.GetAnnotations())
        {
            if (!annotation.Name.StartsWith(RlsAnnotationNames.Prefix,
                                            StringComparison.Ordinal))
                continue;
            // Opt-out-Marker nicht weitergeben, sonst entstehen leere Migrationen
            if (annotation.Name == RlsAnnotationNames.Enabled && annotation.Value is false)
                continue;
            yield return annotation;
        }
    }
}
\end{sourcecode}

Ab hier übernimmt EF: Ändert sich eine Policy, erzeugt der Differ eine
\code{Alter\allowbreak{}Table\allowbreak{}Operation} mit neuen und alten Annotations. In der
generierten Migration sieht das so aus, ganz ohne eigenen Code-Generator:

\begin{sourcecode}[Generierte Migration (\allowbreak{}Ausschnitt)]{csharp}
migrationBuilder.AlterTable(name: "Orders")
    .Annotation("Acme:Rls:Policy:tenant_isolation", "{...neu...}")
    .OldAnnotation("Acme:Rls:Policy:tenant_isolation", "{...alt...}");
\end{sourcecode}

Weil der Annotation-Provider die Annotations nur für \code{design\allowbreak{}Time} weitergibt,
fehlen sie in einem Modell: dem Laufzeitmodell, das \code{context.\allowbreak{}Model}
liefert.

\begin{sourcecode}[Startprüfung eines Plugins (\allowbreak{}FALSCH)]{csharp}
// Hat jede isolierte Tabelle RLS? Gelesen aus dem Laufzeitmodell
var table = context.Model.GetRelationalModel().FindTable("Orders", null);
var enabled = table?[RlsAnnotationNames.Enabled];          // immer null
\end{sourcecode}

\begin{failure}[Was passiert, wenn Plugin-Code das Laufzeitmodell statt des Design-Time-Modells liest?]
  \textbf{Beobachtbar.} Code, der Schema-Annotations des Plugins aus dem relationalen
  Laufzeitmodell liest, findet keine: Eine Prüfung meldet nichts, ein selbst erzeugtes Skript
  enthält keine Policies. Den Differ schützt EF selbst:
  \code{Get\allowbreak{}Differences} mit dem Laufzeitmodell bricht mit dem Hinweis auf
  \code{IDesign\allowbreak{}Time\allowbreak{}Model} ab.

  \textbf{Ursache.} Das Laufzeitmodell ist schlank. Der Annotation-Provider gibt Annotations
  für Tabellen nur zur Design-Zeit weiter, damit sie die Laufzeit nicht belasten.

  \textbf{Lösung.} Für alles, was mit dem Schema zu tun hat,
  \code{context.\allowbreak{}Get\allowbreak{}Service<\allowbreak{}IDesign\allowbreak{}Time\allowbreak{}Model>(\allowbreak{}).\allowbreak{}Model}
  verwenden.

  \textbf{Test:}
  \code{Failure\allowbreak{}Case\allowbreak{}Tests.\allowbreak{}Table\allowbreak{}\_annotations\allowbreak{}\_of\allowbreak{}\_the\allowbreak{}\_runtime\allowbreak{}\_model\allowbreak{}\_…}
  und \code{Differ\allowbreak{}\_rejects\allowbreak{}\_the\allowbreak{}\_runtime\allowbreak{}\_model\allowbreak{}\_…}.
\end{failure}

\section{Der SQL-Generator}

Auch den \code{IMigrations\allowbreak{}Sql\allowbreak{}Generator} gibt es nur einmal. Das
Plugin erbt vom Npgsql-Generator und ergänzt \code{Create\allowbreak{}Table} und
\code{Alter\allowbreak{}Table}:

% source: src/Acme.EntityFrameworkCore.Tenancy.PostgreSQL/Migrations/Internal/TenancyMigrationsSqlGenerator.cs (strings)
\begin{sourcecode}[Migrations/\allowbreak{}Internal/\allowbreak{}Tenancy\allowbreak{}Migrations\allowbreak{}Sql\allowbreak{}Generator.\allowbreak{}cs (\allowbreak{}Auszug)]{csharp}
protected override void Generate(
    CreateTableOperation operation, IModel? model,
    MigrationCommandListBuilder builder, bool terminate = true)
{
    // Unterminiertes SQL kann keine weiteren Statements aufnehmen (wie bei Npgsql)
    if (!terminate && operation.HasRlsConfiguration())
        throw new ArgumentException(
            "RLS kann nicht in unterminiertem CREATE TABLE entstehen.");

    base.Generate(operation, model, builder, terminate);
    if (terminate)
        GenerateRowLevelSecurity(operation.Name, operation.Schema, previous: null,
            operation, builder);
}

protected override void Generate(
    AlterTableOperation operation, IModel? model, MigrationCommandListBuilder builder)
{
    base.Generate(operation, model, builder);
    GenerateRowLevelSecurity(operation.Name, operation.Schema, operation.OldTable,
        operation, builder);
}
\end{sourcecode}

\begin{sourcecode}[Migrations/\allowbreak{}Internal/\allowbreak{}Tenancy\allowbreak{}Migrations\allowbreak{}Sql\allowbreak{}Generator.\allowbreak{}cs (\allowbreak{}gekürzt, vereinfacht)]{csharp}
protected virtual void GenerateRowLevelSecurity(string tableName, string? schema,
    IReadOnlyAnnotatable? previous, IReadOnlyAnnotatable current,
    MigrationCommandListBuilder builder)
{
    var table = Delimit(tableName, schema);
    var oldPolicies = previous?.GetRlsPolicies() ?? new Dictionary<string, string>();
    var newPolicies = current.GetRlsPolicies();

    // 1. Entfernte oder geänderte Policies löschen
    foreach (var (name, json) in oldPolicies)
        if (!newPolicies.TryGetValue(name, out var n) || n != json)
            Statement(builder, $"DROP POLICY {Delimit(name)} ON {table}");

    // 2. RLS-Status (ENABLE/DISABLE, FORCE/NO FORCE)
    if ((previous?.IsRlsEnabled() ?? false) != current.IsRlsEnabled())
        Statement(builder, $"ALTER TABLE {table} "
            + (current.IsRlsEnabled() ? "ENABLE" : "DISABLE") + " ROW LEVEL SECURITY");
    // ... FORCE analog

    // 3. Neue oder geänderte Policies anlegen
    foreach (var (name, json) in newPolicies)
        if (!oldPolicies.TryGetValue(name, out var o) || o != json)
            Statement(builder, CreatePolicySql(table, RlsPolicy.Deserialize(json)));
}

protected virtual void Statement(MigrationCommandListBuilder builder, string sql)
{
    builder.Append(sql).AppendLine(Dependencies.SqlGenerationHelper.StatementTerminator);
    EndStatement(builder);
}
\end{sourcecode}

\begin{datatable}{H{0.60}H{1.40}}
  \textbf{Detail} & \textbf{Begründung} \\
  \tablehead
  DROP + CREATE statt \code{ALTER POLICY} & \code{ALTER POLICY} kann weder \code{FOR} noch
  PERMISSIVE/RESTRICTIVE ändern. DDL ist in PostgreSQL transaktional; es entsteht keine
  sichtbare Lücke. \\
  \tablerule
  \code{End\allowbreak{}Statement(\allowbreak{}builder)} & Beendet den Befehl wie die
  Basisklasse; berücksichtigt dabei Transaktionsunterdrückung. \\
  \tablerule
  Eigene \code{Delimit}-Helfer & Npgsqls \code{Delimit\allowbreak{}Identifier}-Helfer sind
  privat. Bezeichner immer über
  \code{Dependencies.\allowbreak{}Sql\allowbreak{}Generation\allowbreak{}Helper} maskieren,
  nie per String-Verkettung. \\
  \tablerule
  DropTable, RenameTable & Nichts zu tun: PostgreSQL entfernt bzw. verschiebt Policies mit
  der Tabelle. \\
\end{datatable}

% ------------------------------------------------------------
\chapter{Eigene Migrations-Operationen}\label{ch:custom-migration-operations}

Manches gehört nicht ins Modell, aber in Migrationen: Rollen, Rechte, Extensions,
Daten-Backfills. Dafür bietet das Plugin eigene Operationen an. Sie bestehen aus drei
Teilen: der Operationsklasse, einer Extension-Method auf
\code{Migration\allowbreak{}Builder} und der SQL-Erzeugung im Generator.

% source: src/Acme.EntityFrameworkCore.Tenancy.PostgreSQL/Migrations/Operations/TenancyOperations.cs
\begin{sourcecode}[Migrations/\allowbreak{}Operations/\allowbreak{}Tenancy\allowbreak{}Operations.\allowbreak{}cs (\allowbreak{}Auszug)]{csharp}
public class CreateApplicationRoleOperation : MigrationOperation
{
    public required string Name { get; init; }
    public bool CanLogin { get; init; }       // Passwörter gehören NICHT in Migrationen
}
\end{sourcecode}

% source: src/Acme.EntityFrameworkCore.Tenancy.PostgreSQL/Migrations/Operations/TenancyMigrationBuilderExtensions.cs (strings)
\begin{sourcecode}[Migrations/\allowbreak{}Operations/\allowbreak{}Tenancy\allowbreak{}Migration\allowbreak{}Builder\allowbreak{}Extensions.\allowbreak{}cs (\allowbreak{}Auszug)]{csharp}
namespace Microsoft.EntityFrameworkCore.Migrations;      // in Migrationen immer importiert

public static class TenancyMigrationBuilderExtensions
{
    public static OperationBuilder<CreateApplicationRoleOperation> CreateApplicationRole(
        this MigrationBuilder migrationBuilder, string name, bool canLogin = false)
    {
        var operation = new CreateApplicationRoleOperation
        {
            Name = TenancyIdentifiers.ValidateIdentifier(name, "Rollenname"),
            CanLogin = canLogin,
        };
        migrationBuilder.Operations.Add(operation);
        return new OperationBuilder<CreateApplicationRoleOperation>(operation);
    }
}
\end{sourcecode}

Der Generator verteilt Operationen über die zentrale
\code{Generate(\allowbreak{}Migration\allowbreak{}Operation,\allowbreak{}
.\allowbreak{}.\allowbreak{}.\allowbreak{})}-Methode. Unbekannte Typen gehen an die
Basisklasse:

\begin{sourcecode}[Migrations/\allowbreak{}Internal/\allowbreak{}Tenancy\allowbreak{}Migrations\allowbreak{}Sql\allowbreak{}Generator.\allowbreak{}cs (\allowbreak{}Auszug, vereinfacht)]{csharp}
protected override void Generate(
    MigrationOperation operation, IModel? model, MigrationCommandListBuilder builder)
{
    switch (operation)
    {
        case CreateApplicationRoleOperation createRole: Generate(createRole, builder); break;
        case DropApplicationRoleOperation dropRole:     Generate(dropRole, builder); break;
        case GrantTablePrivilegesOperation grant:       Generate(grant, builder); break;
        default: base.Generate(operation, model, builder); break;
    }
}

protected virtual void Generate(CreateApplicationRoleOperation operation,
    MigrationCommandListBuilder builder)
{
    var literal = Literal(operation.Name);           // 'shop_app'  (für den Vergleich)
    var identifier = Delimit(operation.Name);    // shop_app (CREATE ROLE, ggf. in "...")
    var login = operation.CanLogin ? "LOGIN" : "NOLOGIN";

    builder
        .AppendLine("DO $$")
        .AppendLine("BEGIN")
        .AppendLine("    IF NOT EXISTS (SELECT FROM pg_catalog.pg_roles")
        .AppendLine($"                   WHERE rolname = {literal}) THEN")
        .AppendLine($"        CREATE ROLE {identifier} {login} NOSUPERUSER NOBYPASSRLS;")
        .AppendLine("    END IF;")
        .AppendLine("END")
        .Append("$$");
    builder.AppendLine(Dependencies.SqlGenerationHelper.StatementTerminator);
    EndStatement(builder);
}
\end{sourcecode}

\begin{important}[Literal oder Bezeichner]
  Der Rollenname kommt zweimal vor, aber in zwei Formen: als \textbf{String-Literal} im
  Vergleich (\code{'shop\_\allowbreak{}app'}, erzeugt über das String-Type-Mapping) und als
  \textbf{Bezeichner} in \code{CREATE ROLE} (\code{shop\_\allowbreak{}app}, bei Bedarf in
  doppelten Anführungszeichen). Wer beides verwechselt oder per String-Verkettung baut,
  erzeugt ungültiges SQL oder eine Injection-Lücke. Zusätzlich validiert die
  Extension-Method den Namen.
\end{important}

\begin{tip}
  Operationen idempotent formulieren (\code{IF NOT EXISTS}, \code{DROP
  .\allowbreak{}.\allowbreak{}.\allowbreak{} IF EXISTS}). Rollen existieren clusterweit und
  überleben Datenbanken; eine Migration, die beim zweiten Ausführen in einer neuen Datenbank
  scheitert, ist ein häufiger Ärger.
\end{tip}

In einer Migration nutzt man die Operationen wie EFs eigene:

\begin{sourcecode}[Migrations/\allowbreak{}20261021\_\allowbreak{}Initial.\allowbreak{}cs (\allowbreak{}handergänzt)]{csharp}
protected override void Up(MigrationBuilder migrationBuilder)
{
    // VOR den Tabellen (Kapitel „Migrations für Fortgeschrittene“)
    migrationBuilder.CreateApplicationRole("shop_support");
    // ... von EF generierte Operationen (Tabellen, RLS über Annotations) ...
    migrationBuilder.CreateApplicationRole("shop_app", canLogin: true);
    migrationBuilder.GrantTablePrivileges("shop_app");
}

protected override void Down(MigrationBuilder migrationBuilder)
{
    migrationBuilder.DropApplicationRole("shop_app");
    // ...
}
\end{sourcecode}

\begin{note}[Wenn der Differ eigene Operationen erzeugen soll]
  Sollen eigene Operationen vom \emph{Differ} automatisch erzeugt werden, braucht es einen
  abgeleiteten \code{IMigrations\allowbreak{}Model\allowbreak{}Differ} und einen
  \code{ICSharp\allowbreak{}Migration\allowbreak{}Operation\allowbreak{}Generator}, der die
  Operation als C\# schreibt. Beides sind Einzeldienste mit viel interner API. Fast immer
  ist der deklarative Weg über Annotations
  (Kapitel~\ref{ch:declarative-sql-from-annotations}) die bessere Wahl.
\end{note}

% ------------------------------------------------------------
\chapter{Design-Time-Services}\label{ch:design-time-services}

Die Werkzeuge von \code{dotnet ef} (Migrationen anlegen, Snapshots schreiben, Scaffolding,
Compiled Models) verwenden einen eigenen Service Provider mit Design-Time-Diensten. Ein
Plugin erweitert ihn über \code{IDesign\allowbreak{}Time\allowbreak{}Services}.

\section{Automatisch geladen}

Damit Nutzer nichts konfigurieren müssen, verweist ein Assembly-Attribut auf die Klasse. Das
Tooling findet es in jeder Assembly, die das Startprojekt referenziert. Das zweite Argument
beschränkt das Laden auf Contexts mit dem passenden Provider:

% source: src/Acme.EntityFrameworkCore.Tenancy.PostgreSQL/Properties/AssemblyInfo.cs
\begin{sourcecode}[Properties/\allowbreak{}Assembly\allowbreak{}Info.\allowbreak{}cs]{csharp}
[assembly: DesignTimeServicesReference(
    "Acme.EntityFrameworkCore.Tenancy.PostgreSQL.Design.TenancyDesignTimeServices, "
        + "Acme.EntityFrameworkCore.Tenancy.PostgreSQL",
    "Npgsql.EntityFrameworkCore.PostgreSQL")]
\end{sourcecode}

% source: src/Acme.EntityFrameworkCore.Tenancy.PostgreSQL/Design/TenancyDesignTimeServices.cs
\begin{sourcecode}[Design/\allowbreak{}Tenancy\allowbreak{}Design\allowbreak{}Time\allowbreak{}Services.\allowbreak{}cs]{csharp}
public class TenancyDesignTimeServices : IDesignTimeServices
{
    public void ConfigureDesignTimeServices(IServiceCollection services)
    {
        // Alle Registrierungen entfernen (EF-Default, Npgsql), dann unsere setzen:
        // unabhängig von der Reihenfolge, in der das Tooling Services lädt
        services.RemoveAll<IAnnotationCodeGenerator>();
        services.AddSingleton<IAnnotationCodeGenerator, TenancyAnnotationCodeGenerator>();
    }
}
\end{sourcecode}

\section{Lesbare Snapshots: der Annotation-Code-Generator}

Ohne Zutun schreibt EF unbekannte Annotations als
\code{Has\allowbreak{}Annotation(\allowbreak{}.\allowbreak{}.\allowbreak{}.\allowbreak{})}
in den Snapshot, bei Policies mit JSON-Werten. Das funktioniert, ist aber kaum lesbar. Ein
eigener \code{IAnnotation\allowbreak{}Code\allowbreak{}Generator} übersetzt sie in die
Fluent API:

\begin{sourcecode}[Shop\allowbreak{}Context\allowbreak{}Model\allowbreak{}Snapshot.\allowbreak{}cs]{csharp}
// Vorher
b.HasAnnotation("Acme:Rls:Enabled", true);
b.HasAnnotation("Acme:Rls:Policy:tenant_isolation", "{\"Name\":\"tenant_isolation\",...}");

// Nachher
RlsEntityTypeBuilderExtensions.HasRowLevelSecurity(b, true);
RlsEntityTypeBuilderExtensions.HasPolicy(b, "tenant_isolation", p => p
    .Using("\"TenantId\" = nullif(current_setting('app.tenant_id', true), '')::uuid")
    .WithCheck("\"TenantId\" = nullif(current_setting('app.tenant_id', true), '')::uuid"));
\end{sourcecode}

EF Core 10 schreibt Erweiterungsmethoden im Snapshot als statische Aufrufe
(\code{Extensions.\allowbreak{}Methode(\allowbreak{}b,\allowbreak{}
.\allowbreak{}.\allowbreak{}.\allowbreak{})}) statt als
\code{b.\allowbreak{}Methode(\allowbreak{}.\allowbreak{}.\allowbreak{}.\allowbreak{})}. Das
vermeidet Mehrdeutigkeiten, wenn mehrere Provider oder Plugins gleichnamige Methoden
definieren. Tests, die den erzeugten Code prüfen, sollten diese Form erwarten.

\begin{sourcecode}[Design/\allowbreak{}Internal/\allowbreak{}Tenancy\allowbreak{}Annotation\allowbreak{}Code\allowbreak{}Generator.\allowbreak{}cs (\allowbreak{}Auszug, vereinfacht)]{csharp}
public override IReadOnlyList<MethodCallCodeFragment> GenerateFluentApiCalls(
    IEntityType entityType, IDictionary<string, IAnnotation> annotations)
{
    var calls = new List<MethodCallCodeFragment>();

    // VERTRAG: Übersetzte Annotations aus dem Dictionary ENTFERNEN
    if (annotations.TryGetValue(RlsAnnotationNames.Enabled, out var enabled)
        && enabled.Value is true)
    {
        annotations.Remove(RlsAnnotationNames.Enabled);
        var forced = annotations.Remove(RlsAnnotationNames.Forced, out var f)
                     && f.Value is true;
        calls.Add(new MethodCallCodeFragment(HasRowLevelSecurityMethod, forced));
    }

    foreach (var key in annotations.Keys
                 .Where(k => k.StartsWith(RlsAnnotationNames.PolicyPrefix,
                                          StringComparison.Ordinal))
                 .ToList())                         // Kopie: Dictionary ändert sich
    {
        var policy = RlsPolicy.Deserialize((string)annotations[key].Value!);
        annotations.Remove(key);
        calls.Add(new MethodCallCodeFragment(HasPolicyMethod, policy.Name,
            new NestedClosureCodeFragment("p", BuildPolicyChain(policy))));
    }

    calls.AddRange(base.GenerateFluentApiCalls(entityType, annotations));
    return calls;
}
\end{sourcecode}

\begin{datatable}{H{0.70}H{1.30}}
  \textbf{Stolperstelle} & \textbf{Lösung} \\
  \tablehead
  Übersetzte Annotation bleibt im Dictionary & Erscheint doppelt im Snapshot. Immer
  entfernen; was übrig bleibt, schreibt EF als \code{Has\allowbreak{}Annotation}. \\
  \tablerule
  Generische Methode im \code{Method\allowbreak{}Info} & Der Snapshot nutzt den
  nicht-generischen \code{Entity\allowbreak{}Type\allowbreak{}Builder}. Nicht-generische
  Überladung per \code{Get\allowbreak{}Runtime\allowbreak{}Method} mit exakten
  Parametertypen wählen. \\
  \tablerule
  \code{params}-Argument als Array &
  \code{Method\allowbreak{}Call\allowbreak{}Code\allowbreak{}Fragment} nimmt \code{params
  object?[]}. Ein \code{string[]} würde wegen Array-Kovarianz zur Argumentliste: in
  \code{new object?[] \{ roles \}} verpacken. \\
  \tablerule
  Enums im Snapshot & Werden als \code{Policy\allowbreak{}Command.\allowbreak{}Select}
  gerendert. Das Enum liegt deshalb im immer importierten Namespace
  \code{Microsoft.\allowbreak{}Entity\allowbreak{}Framework\allowbreak{}Core}. \\
  \tablerule
  Lambda-Argumente &
  \code{Nested\allowbreak{}Closure\allowbreak{}Code\allowbreak{}Fragment(\allowbreak{}\textquotedbl{}p\textquotedbl{},\allowbreak{}
  kette)} erzeugt \code{p =>
  p.\allowbreak{}Using(\allowbreak{}.\allowbreak{}.\allowbreak{}.\allowbreak{})}; Ketten
  entstehen mit
  \code{fragment.\allowbreak{}Chain(\allowbreak{}.\allowbreak{}.\allowbreak{}.\allowbreak{})}.
  \\
\end{datatable}

Der Generator betrifft nur Snapshot und Designer-Dateien. Die eigentlichen Migrationen
behalten
\code{.\allowbreak{}Annotation(\allowbreak{}.\allowbreak{}.\allowbreak{}.\allowbreak{})} an
\code{Alter\allowbreak{}Table}; das ist der vorgesehene Mechanismus für Änderungen.

% ------------------------------------------------------------
\chapter{Compiled Models, Trimming und NativeAOT}\label{ch:compiled-models-trimming-and-nativeaot}

\code{dotnet ef dbcontext optimize} erzeugt ein \textbf{Compiled Model}: C\#-Code, der das
Laufzeitmodell direkt aufbaut, statt es beim Start über Conventions zu berechnen. Das
verkürzt die Startzeit großer Modelle deutlich und ist Voraussetzung für NativeAOT. Ein
Plugin muss dafür an vier Stellen vorbereitet sein:

\begin{datatable}{H{0.66}H{1.32}H{1.02}}
  \textbf{Baustein} & \textbf{Anforderung} & \textbf{Im Beispiel} \\
  \tablehead
  Type Mappings & Öffentlicher Typ mit öffentlicher statischer Property \code{Default} (für
  NativeAOT genutzt); Konverter/Comparer als gemeinsame Instanzen &
  \code{Tenant\allowbreak{}Id\allowbreak{}Type\allowbreak{}Mapping.\allowbreak{}Default},
  \code{Tenant\allowbreak{}Id\allowbreak{}Converter.\allowbreak{}Instance},
  \code{Tenant\allowbreak{}Id\allowbreak{}Value\allowbreak{}Comparer} \\
  \tablerule
  JSON-Reader/Writer & Statische \code{Instance} und
  \code{Constructor\allowbreak{}Expression} &
  \code{Tenant\allowbreak{}Id\allowbreak{}Json\allowbreak{}Reader\allowbreak{}Writer} \\
  \tablerule
  Value Generators & Konfiguriert mit
  \code{Has\allowbreak{}Value\allowbreak{}Generator\allowbreak{}Factory(\allowbreak{}Type)};
  Factory mit öffentlichem parameterlosem Konstruktor &
  \code{Tenant\allowbreak{}Id\allowbreak{}Value\allowbreak{}Generator\allowbreak{}Factory}
  (Kapitel~\ref{ch:value-generation}) \\
  \tablerule
  Annotations & Einfache Werte übernimmt EF automatisch; ein
  \code{ICSharp\allowbreak{}Runtime\allowbreak{}Annotation\allowbreak{}Code\allowbreak{}Generator}
  filtert Design-Time-Annotations und übersetzt komplexe Werte & Filtert RLS-Policies,
  behält die Laufzeit-Annotation (Kapitel~\ref{ch:property-level-fluent-api-and}) \\
\end{datatable}

Der dritte Punkt wird oft missverstanden. EF übernimmt Annotations mit einfachen Werten
(Strings, Zahlen, bool, Enums) \textbf{automatisch} ins Compiled Model; das zeigt der
Quellcode von
\code{CSharp\allowbreak{}Runtime\allowbreak{}Annotation\allowbreak{}Code\allowbreak{}Generator},
der sie mit \code{Generate\allowbreak{}Simple\allowbreak{}Annotations} schreibt. Ein eigener
Generator ist deshalb nicht nötig, damit Annotations \emph{erhalten} bleiben, sondern um
reine Design-Time-Annotations \emph{herauszufiltern}, die zur Laufzeit niemand liest, und um
Annotations mit komplexen Werten in Code zu übersetzen.
Kapitel~\ref{ch:property-level-fluent-api-and} zeigt beides im Beispiel.

\section{Trimming und AOT}

Alle drei Pakete setzen
\code{<\allowbreak{}Is\allowbreak{}Aot\allowbreak{}Compatible>true<\allowbreak{}/\allowbreak{}Is\allowbreak{}Aot\allowbreak{}Compatible>}.
Der Compiler meldet dann Code, der Reflection ohne Annotationen nutzt. Die
JSON-Serialisierung der Policies verwendet einen Source-Generator statt Reflection:

% source: src/Acme.EntityFrameworkCore.Tenancy.PostgreSQL/Metadata/RlsPolicy.cs
\begin{sourcecode}[Metadata/\allowbreak{}Rls\allowbreak{}Policy.\allowbreak{}cs (\allowbreak{}Auszug)]{csharp}
[JsonSourceGenerationOptions(UseStringEnumConverter = true)]
[JsonSerializable(typeof(RlsPolicy))]
internal partial class RlsJsonContext : JsonSerializerContext;
\end{sourcecode}

Wo Dynamik unvermeidbar ist, wird die Warnung mit Begründung unterdrückt statt übergangen. Die
Isolations-Convention baut den Query-Filter als Expression Tree, was dynamischen Code
braucht. Conventions laufen aber nur beim Modellbau, und NativeAOT setzt ein Compiled Model
voraus, in dem sie nicht mehr laufen; die Unterdrückung nennt genau diesen Grund. Die
Warnungen erscheinen nur bei einem vollständigen Build zuverlässig, inkrementelle Builds
verschweigen sie (Upgrade-Playbook in Kapitel~\ref{ch:versioning-and-compatibility}).

\begin{note}
  Conventions, Snapshots und Migrations laufen nur zur Design-Zeit oder beim Modellbau; mit
  Compiled Model und vorkompilierten Abfragen fallen sie zur Laufzeit weg. Für AOT kritisch
  sind deshalb vor allem Type Mappings, Konverter, Translators (laufen beim Build) und
  Interceptors.
\end{note}

\section{Was als Verbraucher tatsächlich geprüft ist}\label{sec:consumer-test}

\code{Is\allowbreak{}Aot\allowbreak{}Compatible} lässt Analyzer laufen, beweist aber nichts zur
Laufzeit. Auch ein Unit-Test, der Compiled-Model-Code erzeugt und Textstellen prüft
(\code{Compiled\allowbreak{}Model\allowbreak{}Tests}), kompiliert und lädt das Modell nicht.
Deshalb prüft \code{eng/\allowbreak{}consumer-\allowbreak{}test/\allowbreak{}run.\allowbreak{}sh}
das Plugin so, wie eine Anwendung es bekommt, in einem temporären Verzeichnis, das am Ende
gelöscht wird:

\begin{datatable}{H{0.55}H{1.45}}
  \textbf{Schritt} & \textbf{Prüfung} \\
  \tablehead
  Pakete & \code{dotnet pack} der drei Bibliotheken in eine lokale Quelle; ein frisches
  Projekt referenziert die Pakete, nicht die Projekte. Eigener NuGet-Cache, Paketquellen-Zuordnung
  \code{Acme.*} nur aus der lokalen Quelle und ein SHA-512-Vergleich belegen, dass genau die
  eben gebauten Pakete aufgelöst werden. \\
  \tablerule
  Tooling & \code{dotnet-ef} als lokales Tool. Die Design-Time-Dienste werden ohne Zutun
  gefunden. Migration und SQL-Skript enthalten RLS, \code{FORCE}, Policy und Index
  (PostgreSQL) sowie Security Policy und Session-Schlüssel (SQL Server).
  \code{has-\allowbreak{}pending-\allowbreak{}model-\allowbreak{}changes} meldet danach keine
  Änderung. \\
  \tablerule
  Laufzeitmodell & \code{TenantId} wird \code{uuid}, die Laufzeit-Annotation nennt die
  Property, \code{Current\allowbreak{}Tenant(\allowbreak{})} und der Query-Filter werden
  übersetzt. \\
  \tablerule
  Compiled Model & \code{dotnet ef dbcontext optimize}, kompilieren, mit
  \code{Use\allowbreak{}Model} laden, dieselben Prüfungen. \\
  \tablerule
  NativeAOT (\code{-\allowbreak{}-nativeaot}) & Veröffentlichen mit
  \code{Publish\allowbreak{}Aot} und vorkompilierten Abfragen, dann gegen PostgreSQL in Docker
  starten: zwei Mandanten schreiben über die Interceptors, einer liest, RLS liefert nur seine
  Zeilen. Keine Trim- oder AOT-Warnung darf aus einer Acme-Assembly stammen. \\
\end{datatable}

Die CI führt den Ablauf in einem eigenen Job mit dem neuesten stabilen 10.0.x-Patch für das
Tooling aus und gibt die Versionen aus. Der Test hat dabei Dinge gefunden, die die Unit-Tests
nicht zeigen konnten:

\begin{itemize}
  \item \textbf{Query-Filter und Compiled Models schließen sich aus.} EF Core~10 lehnt
    \code{optimize} ab: „Compiled model can't be generated, because query filters are not
    supported.“ Wer den optionalen Filter des Plugins einschaltet, verzichtet auf Compiled
    Model und NativeAOT; der Schutz kommt dann allein von RLS. Der Test erwartet diesen Fehler
    ausdrücklich, damit eine spätere EF-Version, die Filter kann, auffällt.
  \item \textbf{Value Comparer per Reflection.} Das Type Mapping für \code{TenantId} brachte
    keine eigenen Comparer mit. Das Compiled Model verweist dann nur auf
    \code{Tenant\allowbreak{}Id\allowbreak{}Type\allowbreak{}Mapping.\allowbreak{}Default}, und
    EF erzeugt den Comparer beim ersten Tracking mit \code{Make\allowbreak{}Generic\allowbreak{}Method},
    was die NativeAOT-Anwendung abbrechen ließ. Seitdem tragen beide Mappings gemeinsame
    Instanzen aus \code{Tenant\allowbreak{}Id\allowbreak{}Value\allowbreak{}Comparer}, so wie
    Konverter und JSON-Reader.
  \item \textbf{Grenzen von EF selbst.} Vorkompilierte Abfragen und NativeAOT sind in EF
    Core~10 experimentell (Warnung EF9100). \code{Publish\allowbreak{}Aot} darf nicht im
    Projekt stehen, weil EF sonst schon zur Design-Zeit den Modellbau verweigert. Das Paket
    \code{Microsoft.\allowbreak{}Entity\allowbreak{}Framework\allowbreak{}Core.\allowbreak{}Tasks}
    erzeugt Modell und Abfragen beim Veröffentlichen, ein zusätzlich manuell erzeugtes Modell
    stört es. Die Vorkompilierung lehnt \code{o.\allowbreak{}Total > 0} bei einer
    \code{decimal}-Property ab, \code{0m} geht. EF, EF Relational und Npgsql melden selbst
    Trim- und AOT-Warnungen (IL2104, IL3053).
\end{itemize}

Nicht geprüft ist SQL Server unter NativeAOT; das Paket ist dafür markiert, ein Lauf mit
einer echten Datenbank fehlt.


% ------------------------------------------------------------
\chapter{Scaffolding (Reverse Engineering)}\label{ch:scaffolding-reverse-engineering}

\code{dotnet ef dbcontext scaffold} liest ein bestehendes Datenbankschema und erzeugt
Entitäten und einen Context. Ein Plugin, das eigene Schema-Elemente verwaltet, kann sie
dabei ebenfalls erkennen. Der Weg führt über zwei Design-Time-Dienste:

\begin{datatable}{H{0.60}H{1.40}}
  \textbf{Dienst} & \textbf{Aufgabe} \\
  \tablehead
  \code{IDatabase\allowbreak{}Model\allowbreak{}Factory} & Liest das Schema in ein
  \code{Database\allowbreak{}Model}. Ein abgeleiteter Npgsql-Factory könnte
  \code{pg\_\allowbreak{}policies} und
  \code{pg\_\allowbreak{}class.\allowbreak{}relrowsecurity} abfragen und als Annotations an
  \code{Database\allowbreak{}Table} hängen. \\
  \tablerule
  \code{IAnnotation\allowbreak{}Code\allowbreak{}Generator} & Schreibt diese Annotations als
  Fluent API in den generierten \code{On\allowbreak{}Model\allowbreak{}Creating}, mit
  derselben Logik wie beim Snapshot. \\
\end{datatable}

Das Beispiel implementiert beides: Eine eigene Database-Model-Factory liest RLS-Status und
Policies aus der Datenbank, und der Annotation-Code-Generator aus
Kapitel~\ref{ch:design-time-services} schreibt sie in den gescaffoldeten Context.
Kapitel~\ref{ch:scaffolding-in-practice} zeigt die Umsetzung im Detail.
